% Generated by code/scripts/build_full_evaluation_paper.py.
\begin{table}[t]
\centering\small
\caption{Neural and Wikipedia attribute baselines using the same full evaluations, targets, reference sets, layer, prompt correction, gender alternatives and year rule as the contrastive lens.}
\label{tab:full-baselines}
\begin{tabular}{@{}lrrrrr@{}}
\toprule
Model & CEL & NEL & Logit & Diff. logit & Attribute \\
\midrule
OLMo 3 7B & 0.7355 & 0.7650 & 0.2056 & 0.4390 & -0.4016 \\
GPT-J 6B$^{\dagger}$ & 0.4346 & 0.4196 & 0.4669 & 0.2020 & -0.1035 \\
Pythia 6.9B$^{\dagger}$ & 0.5027 & 0.5063 & -0.2047 & 0.1845 & 0.0041 \\
Mistral 7B & 0.6561 & 0.6711 & 0.6581 & 0.5186 & -0.2202 \\
Gemma 2 9B & 0.6889 & 0.6868 & 0.4728 & 0.3436 & -0.5515 \\
Qwen3 8B Base$^{\dagger}$ & 0.6647 & 0.6232 & 0.7908 & 0.5952 & -0.0557 \\
Qwen3.8 27B$^{\dagger}$ & 0.5539 & 0.5717 & 0.3404 & 0.4492 & -0.0099 \\
\bottomrule
\end{tabular}
\par\smallskip\begin{minipage}{\linewidth}\footnotesize CEL: contrastive/differential Jacobian lens; NEL: raw/noncontrastive Jacobian lens. Higher scores mean better target ranks. $^{\dagger}$Leading newline. Attribute overlap is zero when no top-50 predictions survive the published length/stopword filter; no evaluation rows are dropped. These are point estimates, not evidence of significant differences between readouts.\end{minipage}
\end{table}
